\documentclass{article}
\usepackage{amsmath,amssymb,amsthm}
\newtheorem{theorem}{Theorem}
\newtheorem{lemma}{Lemma}
\begin{document}
\title{A062: influence conservation on stopped physical restriction trees}
\author{A062, stream B13}\date{12 October 2026}\maketitle
All entropies are in bits. Let $f:\{-1,1\}^n\to\{-1,1\}$,
$p_f(S)=\widehat f(S)^2$, $H(f)=-\sum_Sp_f(S)\log_2p_f(S)$,
and $I(f)=\sum_S|S|p_f(S)$. These are ordinary influences, not
character-centered influences.

\paragraph{Physical probability space.}
Let $X$ be uniform on the input cube and let $R$ be independent auxiliary
randomness. Include $R$ in $\mathcal F_0$. At time $t$, the next fresh
coordinate $i_t$ (or fresh block $A_t$), and the decision to stop or wait,
are measurable in $\mathcal F_t$. Reveal only the selected coordinates of
$X$. The sigma field $\mathcal F_t$ contains the seed and the revealed
physical history. Let $g_t$ be the actual Boolean restriction of $f$.
Conditionally on this history, all unrevealed coordinates are still
independent fair signs: induction on the reveal history proves this,
since selection observes no unrevealed input. For a block $A_t$, each
assignment has conditional probability $2^{-|A_t|}$. Repeated queries
are redundant and have zero new cost, so omit them. Waiting has zero cost.

\begin{lemma}[One step and fixed blocks]
For a fresh physical block $A$ of a restricted Boolean function $g$,
\[
 I(g)=I_A(g)+2^{-|A|}\sum_z I(g_z),\qquad
 I_A(g)=\sum_{i\in A}\operatorname{Inf}_i(g).
\]
For a singleton $A=\{i\}$, the cost is exactly $\operatorname{Inf}_i(g)$.
\end{lemma}
\begin{proof}
For each $j\notin A$, average the $j$-edge disagreement indicator over
the physical assignment on $A$. This gives
$\operatorname{Inf}_j(g)=2^{-|A|}\sum_z\operatorname{Inf}_j(g_z)$.
Sum over $j\notin A$ and add the influences in $A$.
Equivalently, Fourier Parseval on the revealed block gives the same
identity. If a fixed block is revealed in several stages, possibly in a
history-dependent order, its total expected query cost is $I_A(g)$:
apply the singleton identities and observe that the terminal restriction
has exactly the same law as revealing all of $A$ at once. The realized
path cost need not equal $I_A(g)$.
\end{proof}

\paragraph{Entropy on the same tree.}
Define
\[
 \delta_A(g)=H(g)-2^{-|A|}\sum_zH(g_z).
\]
Let $S\sim p_g$, $q(T)=\Pr(S\setminus A=T)$, and, under a separate
one-node channel experiment, sample $Z$ uniformly on $A$, then sample
$T$ from the Fourier law of $g_Z$. Partial Parseval gives
$q(T)=\mathbb E_Z\widehat {g_Z}(T)^2$. Consequently
\[
 \delta_A(g)=B_A(g)+J_A(g),\quad
 B_A(g)=H(S\cap A\mid S\setminus A),\quad J_A(g)=I(Z;T)\ge0.
\]
Both terms are nonnegative. The physical tree probabilities are not
spectral probabilities; this identity does not choose the next query by
sampling a spectral set.

For a singleton, writing the parent coefficients as $u_T$ and $v_T$ in
the two queried-frequency sectors, the child laws are $(u_T+v_T)^2$ and
$(u_T-v_T)^2$. Their Jensen--Shannon divergence is $J_i$.
Put $m=u^2+v^2$ and $z=2|uv|/m$. In natural units its packet contribution
is $m d(z)$, where
\[
 d(z)=\tfrac12[(1+z)\ln(1+z)+(1-z)\ln(1-z)].
\]
The derivative comparison
$d'(z)=\operatorname{atanh}z\le z/\sqrt{1-z^2}$ proves
$d(z)\le1-\sqrt{1-z^2}$ by integration, including endpoints by
continuity. Thus
\[
 0\le J_i(g)\le\frac{2}{\ln2}\sum_T\min(u_T^2,v_T^2)
 \le\frac{2}{\ln2}\operatorname{Inf}_i(g).                 \tag{1}
\]
The derivative comparison itself follows by differentiating its right
side minus left side: $(1-z^2)^{-3/2}-(1-z^2)^{-1}\ge0$.

\begin{theorem}[Finite-tree optional stopping, including labelled budgets]
Let $\tau$ be a bounded stopping time on the above tree. Let
$c_t=I_{A_t}(g_t)$ and $e_t=\delta_{A_t}(g_t)$, both zero after stopping.
Write $Q_t=\sum_{s<t}c_s$ and $D_t=\sum_{s<t}e_s$. Then
\[
 M_t=I(g_t)+Q_t,\qquad N_t=H(g_t)+D_t
\]
are martingales, and
\[
 I(f)=\mathbb E Q_\tau+\mathbb E I(g_\tau),\qquad
 H(f)=\mathbb E D_\tau+\mathbb E H(g_\tau).                 \tag{2}
\]
In particular, $\mathbb E Q_\tau\le I(f)$.
For each original coordinate $j$, the expected cost of querying $j$ plus
its expected residual influence at the terminal restriction is exactly
$\operatorname{Inf}_j(f)$.
\end{theorem}
\begin{proof}
The one-step identities say
$\mathbb E[I(g_{t+1})\mid\mathcal F_t]=I(g_t)-c_t$ and
$\mathbb E[H(g_{t+1})\mid\mathcal F_t]=H(g_t)-e_t$.
For a stopping time, $\mathbf1_{t<\tau}$ is $\mathcal F_t$-measurable.
Multiply each drift identity by that indicator, take expectations, and
sum over the finite horizon. This proves (2) without invoking any
unbounded optional-stopping theorem.
For a coordinate $j$, keep its residual influence until it is queried,
then replace it by its query cost. Reveals of other coordinates preserve
its conditional expected influence by edge averaging. Its own reveal
adds exactly the quantity removed. This gives a martingale for each
label $j$, and finite summation proves the last assertion.
\end{proof}

If $\pi_v$ is the physical probability of reaching an internal node or
terminal leaf, the same identities read
\[
 I(f)=\sum_{v\text{ internal}}\pi_v c_v+
       \sum_{\ell\text{ terminal}}\pi_\ell I(g_\ell),\quad
 H(f)=\sum_{v\text{ internal}}\pi_v e_v+
       \sum_{\ell\text{ terminal}}\pi_\ell H(g_\ell).       \tag{3}
\]
The weights include all decision randomization as well as the fair
physical branches. A zero-probability node contributes zero.

\paragraph{The payment interface, not a payment theorem.}
For singleton queries, let
$B_{\mathcal T}=\mathbb E\sum_{t<\tau}B_{i_t}(g_t)$ and
$J_{\mathcal T}=\mathbb E\sum_{t<\tau}J_{i_t}(g_t)$.
Then (1)--(2) give
\[
 H(f)=B_{\mathcal T}+J_{\mathcal T}+\mathbb EH(g_\tau),
 \qquad J_{\mathcal T}\le(2/\ln2)\mathbb E Q_\tau.
\]
If a separately proved policy satisfies
$B_{\mathcal T}\le K\mathbb E Q_\tau$ and its actual leaves satisfy
$H(g_\tau)\le C_L I(g_\tau)$ almost surely, then
\[
 H(f)\le \max\{K+2/\ln2,C_L\}I(f).                       \tag{4}
\]
This conclusion is depth independent and charges every label only once.
The policy existence and the bound on $B_{\mathcal T}$ remain open in
the universal setting. A local condition
$\delta_i(g)\le C\operatorname{Inf}_i(g)$ would make
$H(g_t)-C I(g_t)$ a submartingale, not a supermartingale;
equivalently $C I(g_t)-H(g_t)$ is a supermartingale. Neither drift
condition is a consequence of (2).

\paragraph{Fixed dimension with unbounded waiting or randomized stopping.}
On an $n$-cube there are at most $n$ fresh singleton reveals, or $n$ total
coordinates across block reveals. Pathwise $0\le Q_t\le n$,
$0\le I(g_t)\le n$, $0\le H(g_t)\le n$, and
$0\le D_t\le n^2$, since each nonzero entropy decrement is at most $n$.
After absorbing at $\tau$, the restriction eventually stabilizes on each
path, even if $\tau=\infty$: a finite set of actual finite reveal times
has a finite maximum. The terminal restriction means this limiting
restriction. Apply (2) to $\tau\wedge m$ and pass $m\to\infty$ by
bounded convergence. This proves (2) for all such stopping rules,
without a bound on $\mathbb E\tau$. Including independent randomness
in $\mathcal F_0$ justifies arbitrary randomized policies by the same
argument. Anticipative choices using unrevealed inputs are excluded.

\paragraph{Genuinely unbounded processes.}
Suppose an abstract fair-restriction process has nonnegative residual
influence $I_t$, integrable $I_0$, and the stated one-step drift
$\mathbb E[I_{t+1}\mid\mathcal F_t]=I_t-c_t$ with $c_t\ge0$.
For the stopped process, monotone convergence gives
$Q_\infty=\sum_{t<\tau}c_t\in L^1$ and
$\mathbb E Q_\infty\le\mathbb E I_0$.
The nonnegative supermartingale $I_{t\wedge\tau}$ converges almost
surely to $I_*$; Fatou gives only
\[
 \mathbb E Q_\infty+\mathbb E I_*\le\mathbb E I_0.          \tag{5}
\]
Equality follows if the residual influences are uniformly integrable,
for example if dominated by one integrable variable. Equivalently,
given almost-sure convergence, convergence of their expectations is the
needed nonnegative $L^1$ condition. For entropy the same statement
requires finite initial entropy, nonnegative entropy drifts, and a
separate uniform-integrability hypothesis on residual entropy.
Identifying $I_*$ or the entropy limit with a limiting Boolean
restriction is another continuity obligation. No such infinite-input
identification is asserted here.
\end{document}
